% ======================================================================
% Diagrama de flujo — generado automáticamente por generar_diagrama_flujo.py
% NO EDITAR MANUALMENTE; regenerar desde el descriptor JSON.
% ======================================================================
\documentclass[11pt,a4paper]{article}

% ── Codificación, fuentes y lenguaje ───────────────────────────────────────────
\usepackage[utf8]{inputenc}
\usepackage[T1]{fontenc}
\usepackage[default,scale=0.95]{sourcesanspro}
\renewcommand{\familydefault}{\sfdefault}
\usepackage[spanish,es-noshorthands,es-tabla]{babel}

% ── Geometría y espaciado ─────────────────────────────────────────────────────
\usepackage[top=2.2cm, bottom=2.5cm, left=2.2cm, right=2.2cm, headheight=15pt]{geometry}
\usepackage{setspace}
\setstretch{1.18}
\usepackage{parskip}

% ── TikZ (simbolos ISO 5807 / ANSI X3.5) ──────────────────────────────────────
\usepackage{tikz}
\usetikzlibrary{babel, arrows.meta, positioning, shapes.geometric,
                shapes.symbols, decorations.pathmorphing, calc}

% ── Tablas y listas ───────────────────────────────────────────────────────────
\usepackage{booktabs}
\usepackage{tabularx}
\usepackage{array}
\usepackage{enumitem}

% ── Colores, cajas e iconos ───────────────────────────────────────────────────
\usepackage[dvipsnames]{xcolor}
\usepackage{tcolorbox}
\tcbuselibrary{skins, breakable}
\usepackage{fontawesome5}

% ── Secciones con barra de color (infografía) ─────────────────────────────────
\usepackage{titlesec}
\titleformat{\section}
  {\normalfont\LARGE\bfseries\color{azulNoche}}
  {\colorbox{cyanNeon}{\color{fondoOscuro}\thesection}}{0.7em}{}
  [\vspace{2pt}\color{cyanNeon}\rule{\linewidth}{1.8pt}]
\titlespacing*{\section}{0pt}{24pt}{10pt}
\titleformat{\subsection}
  {\normalfont\large\bfseries\color{azulNoche}}
  {\textcolor{cyanNeon}{\thesubsection}}{0.7em}{}
  [\vspace{1pt}\color{grisLinea}\rule{\linewidth}{0.6pt}]
\titlespacing*{\subsection}{0pt}{14pt}{6pt}

% ── Cabeceras y pies de página ────────────────────────────────────────────────
\usepackage{fancyhdr}
\usepackage{lastpage}
\pagestyle{fancy}
\fancyhf{}
\renewcommand{\headrulewidth}{0pt}
\renewcommand{\footrulewidth}{0pt}
\fancyfoot[C]{%
  \begin{minipage}{\linewidth}
    \vspace{2pt}
    \color{cyanNeon}\rule{\linewidth}{1.2pt}\\[2pt]
    \centering\color{grisTexto}\footnotesize
    Página \thepage\ de \pageref{LastPage}
  \end{minipage}%
}

% ── Hiperenlaces ──────────────────────────────────────────────────────────────
\usepackage[colorlinks=true, linkcolor=azulNoche, urlcolor=cyanNeon]{hyperref}
\sloppy
\emergencystretch 3em

% ── Paleta de colores ─────────────────────────────────────────────────────────
\definecolor{fondoOscuro}{HTML}{0D1B2A}
\definecolor{verdeTurquesa}{HTML}{004D40}
\definecolor{azulNoche}{HTML}{1B3A6B}
\definecolor{azulMedio}{HTML}{2A5298}
\definecolor{cyanNeon}{HTML}{00C8E0}
\definecolor{verdeSignal}{HTML}{00B887}
\definecolor{naranjaVivo}{HTML}{FF7043}
\definecolor{rojoAlerta}{HTML}{E53935}
\definecolor{grisPapel}{HTML}{F4F6FA}
\definecolor{grisLinea}{HTML}{C8D0E0}
\definecolor{grisTexto}{HTML}{6B7A99}
\definecolor{amarilloNota}{HTML}{FFB300}

% ── Tarjetas ──────────────────────────────────────────────────────────────────
\tcbset{
  cajaContenido/.style={
    enhanced, breakable,
    arc=5pt, outer arc=5pt,
    colback=grisPapel, colframe=grisLinea,
    boxrule=0.8pt,
    fonttitle=\bfseries\small, coltitle=azulNoche,
    top=6pt, bottom=6pt, left=10pt, right=10pt,
  },
}

% ── Macros ────────────────────────────────────────────────────────────────────
\newcommand{\iconotexto}[2]{\textcolor{cyanNeon}{\faIcon{#1}}~#2}
\newcommand{\iconoBanda}{sitemap}
\newcommand{\bandaTitulo}[3][fondoOscuro]{%
  \noindent
  \begin{tcolorbox}[%
    enhanced, arc=4mm, colback=#1!10!white, colframe=#1, boxrule=0pt,
    borderline west={6pt}{0pt}{cyanNeon},
    top=6pt, bottom=6pt, left=8pt, right=8pt,
  ]
    \noindent
    \begin{minipage}[c]{0.10\linewidth}
      \centering\color{cyanNeon}\faIcon{chalkboard-teacher}
    \end{minipage}%
    \hspace{8pt}%
    \begin{minipage}[c]{0.88\linewidth}
      {\fontsize{20}{24}\selectfont\bfseries\color{white}#2}\par\vspace{5pt}%
      \textcolor{cyanNeon!80!white}{\small\faIcon{angle-right}~#3}%
    \end{minipage}
    \vspace{4pt}
  \end{tcolorbox}%
  \vspace{8pt}%
}


  % ── Estilos de flujo (ISO 5807 / ANSI X3.5) ─────────────────────────────
  \tikzset{
    flujo/.style={
      >=Stealth,
      font=\footnotesize\sffamily,
    },
    terminador/.style={
      draw=azulNoche, fill=azulNoche!8!white, rounded corners=3pt,
      text width=1.6cm, align=center, inner sep=2mm,
      font=\footnotesize\sffamily,
    },
    proceso/.style={
      draw=azulNoche, fill=cyanNeon!10!white,
      text width=1.6cm, align=center, inner sep=2mm,
      font=\footnotesize\sffamily,
    },
    entradasalida/.style={
      draw=azulNoche, fill=verdeSignal!10!white,
      trapezium, trapezium left angle=75, trapezium right angle=105,
      text width=1.8cm, align=center, inner sep=2mm,
      font=\footnotesize\sffamily,
    },
    decision/.style={
      draw=azulNoche, fill=amarilloNota!20!white,
      diamond, aspect=1.6, text width=1.5cm, align=center, inner sep=1pt,
      font=\footnotesize\sffamily,
    },
    almacenamiento/.style={
      draw=azulNoche, fill=grisPapel,
      cylinder, shape border rotate=90, aspect=0.3,
      text width=1.7cm, align=center, inner sep=2mm,
      font=\footnotesize\sffamily,
    },
    documento/.style={
      draw=azulNoche, fill=azulMedio!8!white,
      text width=1.8cm, align=center, inner sep=2mm,
      font=\footnotesize\sffamily,
    },
    nota/.style={
      draw=grisLinea, dashed, fill=grisPapel,
      text width=1.8cm, align=center, inner sep=2mm,
      font=\scriptsize\sffamily,
    },
    etiqueta/.style={
      font=\scriptsize\sffamily, fill=white,
      fill opacity=0.75, text opacity=1, inner sep=1pt,
    },
  }


% ── Cabecera del documento ────────────────────────────────────────────────────
\fancyhead[L]{\small\color{azulNoche}\textbf{ELECTRONICA --- capa 2 (orquestacion)}}
\fancyhead[R]{\small\color{grisTexto}flujo\_charlada.py}
\renewcommand{\headrulewidth}{0pt}

% ─────────────────────────────────────────────────────────────────────────────
\begin{document}

\bandaTitulo{Material de la conferencia}{Generar, verificar por el log y publicar: la charla que usa sus propias tres capas}

\vspace{4pt}
\begin{center}
\begin{tcolorbox}[cajaContenido, title={\faIcon{map-signs}~Leyenda de símbolos---ISO 5807 / ANSI X3.5}]
\small
Óvalo = \textbf{terminador} (inicio/fin) $\cdot$
Rectángulo = \textbf{proceso} (acción determinista) $\cdot$
Rombo = \textbf{decisión} (ramas Sí/No) $\cdot$
Paralelogramo = \textbf{entrada/salida} (lectura/escritura de archivos) $\cdot$
Cilindro = \textbf{almacenamiento online} (estado, snapshots) $\cdot$
Rectángulo con base ondulada = \textbf{documento} (sesiones, logs) $\cdot$
Rectángulo punteado gris = \textbf{nota explicativa} (sin acción).
Flujo: arriba$\rightarrow$abajo, izquierda$\rightarrow$derecha.

\faIcon{tags}~Cada símbolo lleva una \textbf{etiqueta} (T/P/D/E/A/N + número, según su tipo);
su significado se expande en la tabla \emph{Leyenda de etiquetas} bajo cada diagrama.
\end{tcolorbox}
\end{center}

\vspace{6pt}

\section{\iconotexto{cogs}{De la peticion al PDF publicado}}

\begin{center}
\begin{tikzpicture}[flujo, >=Stealth, x=1cm, y=1cm]
  % Fondo decorativo
  \fill[azulNoche!3!white, rounded corners=4pt]
    (-2.3,-0.4) rectangle (5.8,-26.5);
  \node[terminador] (T1) at (0.0,-2.2) {T1};
  \node[proceso] (P1) at (0.0,-4.4) {P1};
  \node[proceso] (P2) at (0.0,-6.6) {P2};
  \node[decision] (D1) at (0.0,-8.8) {D1};
  \node[almacenamiento] (A1) at (0.0,-11.0) {A1};
  \node[decision] (D2) at (0.0,-13.2) {D2};
  \node[proceso] (P3) at (0.0,-15.4) {P3};
  \node[decision] (D3) at (0.0,-17.6) {D3};
  \node[proceso] (P5) at (0.0,-19.8) {P5};
  \node[almacenamiento] (A2) at (0.0,-22.0) {A2};
  \node[terminador] (T2) at (0.0,-24.2) {T2};
  \node[proceso] (P4) at (3.5,-4.4) {P4};
  \node[entradasalida] (E4) at (3.5,-6.6) {E4};
  \node[entradasalida] (E1) at (3.5,-8.8) {E1};
  \node[entradasalida] (E2) at (3.5,-13.2) {E2};
  \node[entradasalida] (E3) at (3.5,-17.6) {E3};

  \draw[->] (T1.south) -- (P1.north);
  \draw[->] (P1.south) -- (P2.north);
  \draw[->] (P1.east) -- (P4.west) node[etiqueta, above, pos=0.5]{--verificar-cifras};
  \draw[->] (P4.south) -- (E4.north);
  \draw[->] (P2.south) -- (D1.north);
  \draw[->] (D1.south) -- (A1.north) node[etiqueta, left, pos=0.5]{Si};
  \draw[->] (D1.east) -- (E1.west) node[etiqueta, above, pos=0.5]{No};
  \draw[->] (A1.south) -- (D2.north);
  \draw[->] (D2.south) -- (P3.north) node[etiqueta, left, pos=0.5]{Si};
  \draw[->] (D2.east) -- (E2.west) node[etiqueta, above, pos=0.5]{No};
  \draw[->] (P3.south) -- (D3.north);
  \draw[->] (D3.south) -- (P5.north) node[etiqueta, left, pos=0.5]{Si};
  \draw[->] (D3.east) -- (E3.west) node[etiqueta, above, pos=0.5]{No};
  \draw[->] (P5.south) -- (A2.north);
  \draw[->] (A2.south) -- (T2.north);
\end{tikzpicture}
\end{center}

\vspace{6pt}
\begin{tcolorbox}[cajaContenido, title={\faIcon{table}~Leyenda de etiquetas}]
\small
\begin{tabularx}{\linewidth}{@{}>{\centering\arraybackslash}p{1.4cm}X@{}}
\toprule
\textbf{Etiqueta} & \textbf{Significado} \\ \midrule
A1 & Leer el informe .tmp/charlada/deck.json: las cifras medidas y, por pieza, errores, overfull y ruta del PDF. \\
A2 & Escribir la vista run\_state\_<run\_id>.json con los pasos y sus codigos. El log append-only no se toca. \\
D1 & La capa 3 salio con codigo 0 y dejo el informe deck.json? \\
D2 & El LOG de cada pieza esta limpio? No se pregunta si el PDF existe: en nonstopmode el PDF existe aunque el documento este roto. \\
D3 & Las cifras del informe coinciden con el disco? \\
E1 & Salida 3: fallo de la capa 3 o informe ausente. No hay nada que publicar. \\
E2 & Salida 3: compilación sucia. No se publica NADA, ni siquiera el material que si compilo. \\
E3 & Salida 1: las cifras no cuadran. Es una inconsistencia del generador, no del entorno. \\
E4 & Salida 0 o 1: cifra confirmada ejecutando 585 aserciones. No se genera material. \\
P1 & Acotar y sanear las entradas: --titulo, --autor, --fecha, --minutos y --salida-dir. Todo se escapa UNA vez aqui; decidir tambien si la peticion es solo --verificar-cifras. \\
P2 & Delegar en la capa 3: generador\_charlada\_latex.py mide el disco, redacta los .tex y compila en .tmp/charlada/. La capa 2 no mide ni compila. \\
P3 & Contrastar las cifras del informe con lo que hay en disco ahora mismo (directivas y scripts). Una cifra no reproducible no se publica. \\
P4 & Variante --verificar-cifras: reejecutar los 6 ficheros de test y sumar sus aserciones. La cifra citada en el atril se vuelve a obtener, no se copia. \\
P5 & Publicar los PDF verificados en la carpeta de entrega (docs/AGENTE\_IA/). Es el UNICO punto donde se copia algo. \\
T1 & Se pide el material de la conferencia: 30 min, audiencia de docentes, eje transferible. \\
T2 & Material verificado y publicado; alerta de exito. \\
\bottomrule
\end{tabularx}
\end{tcolorbox}


\section{\iconotexto{sticky-note}{Notas}}
\begin{itemize}
\item Un PDF existente NO certifica un documento: en modo nonstopmode el PDF se produce aunque este roto. La unica pregunta fiable es si su log tiene lineas '!'.
\item Las tres verificaciones (log limpio, cifras reproducibles, informe presente) las hace la capa 2, no la capa 3: el generador no puede juzgarse a si mismo.
\item El flujo no publica hasta que los dos documentos compilan limpios: si uno falla, no se copia ni el otro.
\item La cifra citada en el atril (585 aserciones) se verifica reejecutando los tests, no copiandola de un documento.
\end{itemize}

% ─────────────────────────────────────────────────────────────────────────────
\end{document}
